\chapter{Conclusion et perspectives}
\label{ch:conc}

\section{Bilan}

Le projet livre une application Web \textbf{complète} répondant aux exigences du module : Blazor, EF Core Code First avec migrations, architecture en couches, CRUD et recherche LINQ, tableau de bord avec graphiques, validation par annotations, code asynchrone. L'extension \textbf{data science} renforce la dimension analytique (risque, tendance, backtest) et supporte la partie « créativité » du barème.

\section{Pistes d'évolution}

\begin{itemize}[leftmargin=*]
  \item \textbf{Authentification} (ASP.NET Core Identity) et portefeuilles utilisateur isolés.
  \item \textbf{Tests unitaires} sur \texttt{PositionCostHelper}, règles de transaction et calculs d'analytique.
  \item \textbf{CI/CD} (GitHub Actions) : build + tests + publication.
  \item \textbf{API REST} optionnelle pour alimenter une mobile app ou un tableau externe.
  \item \textbf{Coût moyen} plus avancé (FIFO réel) et frais de transaction dans le modèle.
  \item Base \textbf{PostgreSQL} ou \textbf{SQL Server} pour un déploiement multi-utilisateurs.
\end{itemize}

\section{Remerciements / répartition travail d'équipe}

\textit{[À compléter : rôles des membres, répartition des tâches, difficultés rencontrées.]}
